% !TEX root = ../main.tex
\chapter{Estat de l'art i context tecnològic}\label{ch:estat_art}
\fancyhead[RE]{\slshape \textsf{Capítol \thechapter.\ Estat de l'art}}
\fancyhead[LO]{\slshape \textsf{\nouppercase \rightmark}}

Aquest capítol situa el treball en el context de la compressió d'imatges
d'observació de la Terra. Primer es presenten els requisits generals de la
compressió multiespectral i els estàndards espacials CCSDS, amb especial
atenció a CCSDS~122 perquè és el referent clàssic més proper a SORTENY.
Després s'introdueixen els còdecs neuronals, el model SORTENY, el concepte de
qualitat fixa i l'ús de regions d'interès. El capítol es tanca amb les
restriccions d'execució a bord i amb el buit pràctic que motiva la contribució
d'aquesta memòria.

\section{Criteri de revisió bibliogràfica}\label{sec:cerca}

La cerca d'informació s'ha organitzat al voltant de quatre blocs: compressió
espacial clàssica, compressió apresa, control de qualitat local i processament
a bord. S'han prioritzat fonts primàries o oficials: recomanacions CCSDS,
normes ISO/IEC i ECSS, articles revisats per parells, documentació de missions
i documentació tècnica dels fabricants. Quan s'han utilitzat fonts de context,
com fitxes de hardware o informació institucional de missions, no s'han emprat
per derivar xifres experimentals pròpies.

La revisió no pretén enumerar totes les variants existents. Cada referència
s'ha inclòs perquè aporta una definició, una alternativa tècnica, un mètode
comparable o una restricció de disseny. Les xifres experimentals pròpies, en
canvi, no depenen de la bibliografia: es vinculen a checkpoints i manifests del
projecte. Aquesta separació evita confondre resultats de tercers amb resultats
generats durant aquesta memòria. El treball previ propi s'utilitza com a
antecedent publicat en línia i es cita explícitament quan se'n reutilitza una
idea, una figura o una comparació \cite{anez2026treball}.

\section{Compressió d'imatges d'observació de la Terra}\label{sec:compressio_remota}

Una imatge multiespectral es pot escriure com un tensor
$\mathbf{x}\in\mathbb{R}^{B\times H\times W}$, on $B$ és el nombre de bandes
espectrals, $H$ l'alçada i $W$ l'amplada. La compressió intenta reduir la mida
del fitxer aprofitant dues redundàncies principals: la redundància espacial
entre píxels propers i la redundància espectral entre bandes correlacionades.
En imatges hiperespectrals, on el nombre de bandes pot ser molt més gran, la
redundància espectral és encara més important \cite{bioucas2013hyperspectral,
guanter2015enmap}.

Hi ha tres règims habituals. La compressió sense pèrdues reconstrueix
exactament les mostres originals, però l'estalvi queda limitat per l'entropia
de la font \cite{cover2006information,sayood2017compression}. La compressió
gairebé sense pèrdues introdueix un error acotat. La compressió amb pèrdues
accepta una distorsió més gran a canvi d'una reducció de dades superior. En
observació de la Terra no hi ha una resposta universal: depèn de la missió, del
sensor, de la calibració i de la tasca científica posterior
\cite{blanes2014tutorial,gonzalez2025robustness}.

Per comparar resultats en imatges multiespectrals, aquesta memòria expressa la
taxa en bits per píxel i per banda:

\begin{equation}
R_{\mathrm{bps}} =
\frac{8\,S_{\mathrm{fitxer}}}{HWB},
\label{eq:bps_intro}
\end{equation}

on $S_{\mathrm{fitxer}}$ és la mida comprimida en bytes. Per tant, \emph{bps}
no vol dir bits per segon en aquest document, sinó bits per mostra espacial i
per banda espectral. Quan es parla de velocitat temporal, s'utilitza
\emph{throughput} o temps d'execució.

JPEG~2000 és una referència general de compressió d'imatge fixa amb i sense
pèrdues \cite{iso15444}. Tot i això, les missions espacials solen prendre com
a punt de partida les recomanacions CCSDS, perquè estan pensades per
interoperabilitat, implementació embarcada i fluxos de dades d'instruments.

\section{Estàndards CCSDS}\label{sec:ccsds}

El Consultative Committee for Space Data Systems (CCSDS) defineix
recomanacions utilitzades en missions espacials. En compressió d'imatges i
dades multiespectrals, tres famílies són especialment rellevants: CCSDS~121,
CCSDS~122 i CCSDS~123.

\subsection{CCSDS 121 i CCSDS 123 com a context}

CCSDS~121 defineix compressió sense pèrdues per a dades d'instruments. Combina
un preprocessament opcional amb codis adaptatius, com els codis Rice, i prioritza
descodificació exacta i baixa complexitat \cite{ccsds121}. És molt útil com a
etapa de codificació per entropia o com a referència de robustesa, però no
defineix per si mateix una transformada espacial amb pèrdues.

CCSDS~123 està orientat a imatges multiespectrals i hiperespectrals. Es basa en
predicció: estima el valor d'una mostra a partir del veïnatge espacial i
espectral, i després codifica el residu \cite{ccsds123}. La versió B-2 permet
funcionament sense pèrdues o gairebé sense pèrdues i diferents codificadors per
entropia \cite{blanes2019parameter,kiely2018new}. Existeixen implementacions
paral·leles i hardware accelerat \cite{bascones2017parallel,barrios2021hls,
johansen2018fpga}, i arquitectures com SHyLoC mostren la maduresa de la línia
CCSDS en entorns espacials \cite{barrios2020shyloc}.

Tot i aquesta rellevància, CCSDS~123 no és el referent principal per situar
aquesta memòria, perquè SORTENY és un compressor amb pèrdues basat en
transformades apreses, no un predictor near-lossless. Per això el punt de
referència conceptual més proper és CCSDS~122.

\subsection{CCSDS 122 com a referent clàssic}

CCSDS~122 és una recomanació d'imatge basada en transformada. La seva estructura
general es pot veure a la Figura~\ref{fig:ccsds122_pipeline}: primer s'aplica
una transformada discreta d'ondeletes (DWT) i després un codificador de plans de
bits (BPE) genera el flux codificat \cite{ccsds122}.

\begin{figure}[H]
\centering
\includegraphics[width=0.72\textwidth]{ch2/ccsds122_pipeline.png}
\caption{Estructura general de CCSDS~122: transformada DWT seguida d'un codificador de plans de bits.}
\label{fig:ccsds122_pipeline}
\figsource{Adaptat del treball previ propi \cite{anez2026treball}.}
\end{figure}

La DWT separa la imatge en subbandes de baixa i alta freqüència. Aplicada de
forma separable, primer filtra files i després columnes, generant components
com $LL$, $LH$, $HL$ i $HH$. CCSDS~122 pot utilitzar filtres 5/3 enters per a
compressió reversible o filtres 9/7 biortogonals per a compressió amb pèrdues.
La transformada es pot aplicar recursivament sobre la subbanda $LL$, formant
una piràmide multiresolució com la de la Figura~\ref{fig:ccsds122_dwt}.

\begin{figure}[H]
\centering
\includegraphics[width=0.82\textwidth]{ch2/ccsds122_dwt_multilevel.png}
\caption{Descomposició multiresolució amb DWT. La subbanda $LL$ concentra la informació de baixa freqüència i es pot tornar a descompondre.}
\label{fig:ccsds122_dwt}
\figsource{Adaptat del treball previ propi \cite{anez2026treball}.}
\end{figure}

El BPE codifica els coeficients transformats per plans de bits. Els bits més
significatius es codifiquen abans que els menys significatius, de manera que el
flux pot truncar-se i encara permet reconstruir una imatge completa amb menys
qualitat. A més, l'estructura d'arbres entre coeficients pare i fills en les
subbandes wavelet permet explotar que, si un coeficient de baixa freqüència és
poc significatiu, sovint també ho són els seus descendents
\cite{poupat2013cwicom}. Aquesta relació es resumeix a la
Figura~\ref{fig:ccsds122_family}.

\begin{figure}[H]
\centering
\includegraphics[width=0.55\textwidth]{ch2/ccsds122_wavelet_family.png}
\caption{Exemple de relació familiar entre coeficients wavelet utilitzada pel BPE.}
\label{fig:ccsds122_family}
\figsource{Adaptat del treball previ propi \cite{anez2026treball}.}
\end{figure}

CCSDS~122 és, per tant, un referent exigent: és madur, interpretable, ràpid i
orientat a implementacions embarcades. En el treball previ propi es va utilitzar
com a comparador experimental per situar el cost i la qualitat de SORTENY C
\cite{anez2026treball}. En aquesta memòria no es repeteix aquella comparativa:
CCSDS~122 es manté com a referència tecnològica de còdec espacial clàssic,
mentre que els experiments se centren en una pregunta diferent: com controlar
espacialment la qualitat dins de SORTENY i quin cost té fer-ho en C.

\section{Compressió apresa i SORTENY}\label{sec:learned}

\subsection{Autoencoders i models taxa--distorsió}

Els còdecs apresos parteixen d'una idea similar a la dels còdecs clàssics:
transformar la imatge a un domini on sigui més fàcil eliminar redundància,
quantitzar i reconstruir. La diferència és que la transformada no es defineix
manualment, com la DWT de CCSDS~122, sinó que s'aprèn a partir de dades. En la
forma bàsica, una transformada d'anàlisi $g_a$ produeix latents
$\mathbf{y}$, aquests latents es quantitzen, i una transformada de síntesi
$g_s$ reconstrueix la imatge:

\begin{equation}
\mathbf{y}=g_a(\mathbf{x}), \qquad
\hat{\mathbf{y}}=Q(\mathbf{y}), \qquad
\hat{\mathbf{x}}=g_s(\hat{\mathbf{y}}).
\end{equation}

Els latents $\mathbf{y}$ són, per tant, la representació comprimible de la
imatge. Si la xarxa aprèn una bona transformada, molts valors latents són petits
o previsibles, i el codificador per entropia pot representar-los amb menys
bits. Per això el còdec no es redueix només a una xarxa neuronal: també cal un
model probabilístic dels latents i una etapa de codificació. Durant
l'entrenament es minimitza una funció taxa--distorsió:

\begin{equation}
\mathcal{L}=R(\hat{\mathbf{y}})+\lambda D(\mathbf{x},\hat{\mathbf{x}}),
\label{eq:rd_loss}
\end{equation}

on $R$ aproxima la taxa, $D$ mesura la distorsió i $\lambda$ regula el
compromís entre totes dues. En termes intuïtius, $R$ penalitza que els latents
siguin difícils de codificar, mentre que $D$ penalitza perdre informació visual
o radiomètrica. Els primers autoencoders end-to-end van establir aquest esquema
\cite{balle2017end,theis2017lossy}.

La qualitat del terme $R$ depèn de com s'estimi la probabilitat dels latents.
Els hiperpriors introdueixen informació lateral per predir millor aquesta
distribució i millorar el codificador aritmètic
\cite{balle2018hyperprior,minnen2018joint}. Els models de context van un pas
més enllà: estimen la probabilitat d'un latent aprofitant informació d'altres
latents ja coneguts \cite{cheng2020learned,he2022elic}. Això pot reduir la taxa,
però sovint introdueix dependències seqüencials i més latència. Per a una
implementació embarcada, aquesta millora no és automàticament acceptable si
augmenta massa memòria o temps.

La segona idea important per a SORTENY és el funcionament multiràtio. En un
còdec après convencional, cada punt taxa--distorsió pot requerir un model
entrenat per separat. Dumas \etal\ van mostrar que una mateixa transformada
podia operar en diversos punts variant el pas de quantització durant la
inferència \cite{dumas2018quantization}. SORTENY resol el problema amb una
estratègia pròpia, basada en un modulador condicionat per $\lambda$, però la
motivació és la mateixa: desacoblar els pesos entrenats del punt concret de
compressió. Aquesta és la porta que després permet assignar qualitats locals.

Eines com CompressAI faciliten la reproducció científica de còdecs neuronals
\cite{begaint2020compressai}. Tot i això, una implementació per a processament
a bord no pot dependre només d'un framework generalista: ha de controlar
memòria, tipus de dada, paral·lelisme, fitxers de pesos i temps d'execució.

\subsection{Compressió neuronal en teledetecció}

L'aprenentatge profund s'ha aplicat àmpliament a dades de teledetecció
\cite{zhu2017deep}. En compressió remota, el repte principal és que la imatge
no és només espacial. Una fotografia RGB convencional té tres canals, mentre que
una escena multiespectral conté bandes que responen a fenòmens físics diferents:
vegetació, aigua, sòl, núvols o absorcions atmosfèriques. Per això un compressor
remot ha de compactar textures i vores, però també la relació entre bandes.

Les propostes neuronals aborden aquest repte de diferents maneres. Alguns
treballs utilitzen convolucions en diverses direccions per capturar dependències
espacials i espectrals sense tractar les bandes com canals independents
\cite{kong2022spectral}. Altres incorporen informació de vores per preservar
detalls estructurals que podrien perdre's amb una compressió massa agressiva
\cite{guo2023edge}. També hi ha treballs que milloren el model estadístic amb
priors globals o models de context de la imatge completa \cite{zhang2024global}.
Tots aquests mecanismes apunten al mateix objectiu: que els latents siguin més
previsibles i, per tant, més compressibles.

El problema és que cada millora del model pot tenir un cost. Els mecanismes
d'atenció, per exemple, aprenen a donar més pes a certes posicions, bandes o
canals latents, però afegeixen operacions i memòria. Els models de context poden
reduir bits, però introdueixen dependències. En processament a bord no n'hi ha
prou amb millorar la corba taxa--distorsió en un ordinador de laboratori: cal
que la inferència sigui viable en CPU, amb poca RAM i sense dependències pesants.

Per això és especialment rellevant la línia de models reduïts per a execució a
bord: menys canals, menys capes, operacions més simples o combinacions de
compressió i denoising \cite{oliveira2021reduced,oliveira2022denoising,
mijares2023reduced}. SORTENY s'inscriu en aquest context. Manté la idea d'una
transformada apresa multiespectral, però incorpora una modulació que permet
canviar el punt de qualitat sense substituir els pesos del model.

L'interès institucional també és creixent. Phi-Sat-2 inclou aplicacions d'IA a
bord, com detecció de núvols i compressió profunda \cite{guerrisi2023phisat,
esa2024phisat}. Això no elimina els requisits de verificació, throughput i
robustesa, però mostra que el processament neuronal embarcat ja és una línia
operativa real.

\subsection{SORTENY}

SORTENY (\emph{Spectral Orthogonal Transform Encoder}) és un compressor
multiespectral basat en autoencoder \cite{mijares2025sorteny}. La Figura
\ref{fig:sorteny_architecture} mostra la seva arquitectura general. El model
combina tres idees: una transformada espectral apresa, una transformada espacial
convolucional i una modulació de qualitat.

\begin{figure}[H]
\centering
\includegraphics[width=0.94\textwidth]{ch2/sorteny_architecture.png}
\caption{Arquitectura conceptual de SORTENY: transformada espectral, transformada espacial, modulació de qualitat, quantització i síntesi.}
\label{fig:sorteny_architecture}
\figsource{Adaptat del treball previ propi \cite{anez2026treball}.}
\end{figure}

La transformada espectral redueix la redundància entre bandes projectant cada
píxel multiespectral a un espai de característiques après. Després, la
transformada espacial utilitza convolucions amb \emph{stride} i normalitzacions
GDN per compactar la informació en un tensor latent més petit. Finalment, una
subxarxa de modulació rep el paràmetre $\lambda$ i genera factors que escalen
els latents abans de la quantització. Un $\lambda$ més alt conserva més
informació i sol requerir més bits; un $\lambda$ més baix força més valors cap
a zero i augmenta la compressió.

Aquesta modulació és la propietat clau per a aquesta memòria. Si el model pot
canviar de qualitat amb un paràmetre, llavors es pot assignar un paràmetre de
qualitat diferent a cada regió de la imatge. El treball fixed-quality de
Mijares i Verdú \etal\ estudia justament com estimar el paràmetre necessari per
assolir un MSE objectiu local \cite{mijares2024fixed}. Aquesta base científica
permet passar d'una compressió uniforme a una compressió amb control espacial
de qualitat.

\section{Qualitat fixa, taxa fixa i regions d'interès}\label{sec:fixed_quality_art}

En compressió és important distingir taxa fixa i qualitat fixa. En un sistema
de taxa fixa, l'usuari imposa el volum màxim de dades; per exemple, ``aquesta
imatge ha d'ocupar com a màxim 2 MB''. El còdec intenta adaptar-se a aquest
límit i la qualitat final és la conseqüència. En un sistema de qualitat fixa,
l'usuari imposa una distorsió objectiu; per exemple, ``aquesta regió ha de
mantenir un PSNR mínim''. En aquest cas, la mida final del fitxer és la
conseqüència.

La qualitat fixa és atractiva en observació de la Terra perquè la utilitat de
la imatge no depèn només de la mida del fitxer, sinó de si les zones
científicament rellevants conserven prou informació. A diferència del
near-lossless, que sol limitar l'error màxim per mostra, el fixed-quality pot
formular objectius de MSE o PSNR sobre blocs o regions
\cite{mijares2024fixed}.

La regió d'interès (ROI) introdueix una pregunta addicional: no només cal
decidir quanta qualitat es vol, sinó on s'ha de preservar. Una imatge pot tenir
bon PSNR global i, alhora, degradar just la zona que l'usuari vol analitzar.
Els mètodes \emph{content-weighted} ja han explorat assignacions espacials de
qualitat en còdecs apresos \cite{li2018content}. En una plataforma embarcada,
però, la decisió espacial ha de ser barata, transparent i verificable.

Els índexs espectrals ofereixen una via simple per definir presets automàtics.
La forma habitual és comparar dues bandes amb un quocient normalitzat:

\begin{equation}
I(A,B)=\frac{A-B}{A+B}.
\label{eq:spectral_index_intro}
\end{equation}

Aquesta normalització redueix l'efecte de l'escala radiomètrica i permet
interpretar contrastos físics de manera aproximada. L'exemple més conegut és
l'NDVI:

\begin{equation}
\mathrm{NDVI}=\frac{\mathrm{NIR}-\mathrm{Red}}
{\mathrm{NIR}+\mathrm{Red}},
\end{equation}

on NIR és l'infraroig proper. Valors positius elevats solen indicar vegetació
activa, perquè la vegetació reflecteix fortament al NIR i absorbeix més al
vermell. Valors propers a zero poden correspondre a sòl nu o cobertes urbanes,
i valors negatius sovint apareixen en aigua, ombres o superfícies amb baixa
resposta NIR. Aquesta lectura és útil, però no és una classificació universal
\cite{rouse1974ndvi}.

Altres índexs segueixen la mateixa filosofia. L'NDWI contrasta verd i infraroig
proper per delinear aigua \cite{mcfeeters1996ndwi}; el GNDVI substitueix el
vermell pel verd per obtenir una variant sensible a vegetació
\cite{gitelson1996gndvi}; i l'NDCI es va proposar per estimar clorofil·la en
aigües productives \cite{mishra2012ndci}. El llindar depèn del sensor, de
l'atmosfera, de la reflectància, de la resolució i de l'escena. Per això
aquesta memòria parla de \emph{presets operatius}: regles interpretables que
seleccionen blocs candidats a rebre més qualitat, no etiquetes físiques
perfectes.

\section{Processament a bord i hardware de baix recurs}\label{sec:onboard_art}

La compressió a bord es justifica perquè l'enllaç de baixada, la memòria i el
temps disponible en un satèl·lit són limitats \cite{selva2012cubesats,
nasa2017cubesat}. Les missions petites i els CubeSats han impulsat l'ús de
components comercials i processadors de baix consum, però també introdueixen
restriccions de potència, temperatura, radiació i verificació
\cite{denis2017disruptions}.

La Raspberry Pi 3B+ utilitzada en aquesta memòria integra quatre nuclis
Cortex-A53 i 1~GB de memòria \cite{raspberry2018brief}. No és hardware
qualificat per a vol i pot patir problemes que no apareixen igual en un
ordinador terrestre, com alimentació insuficient, temperatura o susceptibilitat
a radiació \cite{guertin2021raspberry}. Tot i això, el Cortex-A53 i components
COTS han estat estudiats per a processament espacial \cite{schwaller2017a53},
i el C3SatP embarcat a Menut mostra una arquitectura COTS propera a aquest
rang de recursos \cite{serra2024c3satp}. Per tant, la Raspberry no certifica
vol, però sí és una plataforma raonable per mesurar ordres de magnitud de
temps, memòria i comportament tèrmic.

La implementació d'aquest treball utilitza C i OpenMP. OpenMP permet expressar
paral·lelisme de memòria compartida, però exigeix verificar que no s'introdueixen
carreres ni canvis d'ordre aritmètic \cite{openmp2021}. Per això les
optimitzacions posteriors es validen contra una baseline C abans d'acceptar-les.

Des del punt de vista d'assegurament del producte, aquesta memòria no
desenvolupa software de vol certificat. Les normes ECSS-E-ST-40C i
ECSS-Q-ST-80C indiquen què caldria per a un procés espacial complet:
requisits traçables, control de configuració, validació, verificació, gestió de
fallades i assegurament del producte software \cite{ecss40,ecss80}. Aquí es
cobreixen traçabilitat experimental, tests, manifests i reproducció, però no
qualificació de components, tolerància a fallades ni certificació.

\section{Síntesi del buit tecnològic}\label{sec:gap}

Aquest capítol mostra que les peces existeixen, però normalment apareixen
separades. CCSDS aporta estàndards madurs i eficients; els còdecs neuronals
aporten transformades apreses; el fixed-quality aporta control de distorsió;
els índexs espectrals aporten regles simples per seleccionar zones; i els
treballs d'execució a bord mostren la necessitat de controlar memòria i temps.

El buit que aborda aquesta memòria és la integració: convertir una transformada
apresa multiespectral en una ruta operativa que decideixi automàticament quines
zones protegir, assigni qualitat local, s'executi en C sense TensorFlow durant
l'operació i permeti mesurar tant la taxa codificada com el cost computacional
en hardware limitat.

La novetat del treball no és proposar un nou índex espectral ni un nou
autoencoder. La contribució és construir i validar una arquitectura que uneix
control espacial de qualitat, implementació C, presets automàtics, mesura de
bitrate codificat i avaluació en una plataforma de baix recurs. Els capítols
següents defineixen els fonaments tècnics necessaris i després descriuen la
implementació i els experiments.
